\section{Introduction}

\subsection{Motivation}
Can an AI research agent, operating with modern LLM reasoning plus tool access
and modest HPC, faithfully replicate published computational science? The
question matters for three reasons: (i) reproducibility is the epistemic
backbone of the field, yet wet- and dry-lab replication is expensive;
(ii) agentic AI systems are now capable of non-trivial end-to-end pipelines,
so their replication failure modes and strengths need empirical characterization;
and (iii) a working replication pipeline opens the door to automated follow-on
research --- hypothesis generation, robustness probes, and scaling studies.

This report documents our first sixty attempts, including a second wave
of fifteen replications completed in a 5-day sprint (2026-04-26 to
2026-04-30) via parallel subagents, a third wave of ten
biology/bioinformatics replications (2026-05-05) leveraging the BV-BRC
genomic database API, and a fourth wave of five PDE/PINN replications
(2026-05-07) probing the reproducibility of machine-learning approaches
to partial differential equations.

\subsection{Scope and source}
We drew papers from \href{https://www.osti.gov}{OSTI.gov}, filtered for
Argonne-authored computational work with enough method detail to be
reproducible, plus select community benchmarks in PDE and PINN methods.
The cohort spans $\sim 20$+ domains
(Table~\ref{tab:domain}): astrophysics/cosmology, nuclear physics, condensed
matter, materials, computational chemistry, computational biology, climate
science, fusion energy, NLP/domain language models, quantum
information, quantum optics, power systems, power electronics, combinatorics,
graph algorithms, control theory, scientific ML, PDE/PINN methods, dark matter direct detection,
combustion LES, and others. Publication years range from 2002 to 2026.

\subsection{Methodology}
Each replication was executed by an agentic pipeline consisting of:
(a) \textbf{paper ingestion} (PDF $\rightarrow$ structured plan),
(b) \textbf{environment setup} (uv/pip/conda, container images, data
acquisition),
(c) \textbf{implementation} (write or adapt code, usually in
Python/Fortran/C++),
(d) \textbf{execution} on the appropriate compute tier
(CherryRd Mac Studio for small jobs; uicgpu 8$\times$A100 for mid-scale
DL training; Aurora or Polaris for HPC ensembles), and
(e) \textbf{scoring and reporting} against the published numbers.
The model backbone used was Anthropic Claude Opus 4 via a free Argo proxy at
Argonne (Wave~1 used GPT-family models; Wave~2 switched to Claude Opus as the
default). Wave~2 introduced \textbf{subagent parallelism}: the main agent
spawned independent subagent sessions, each responsible for a single paper,
running concurrently on the same compute fabric. A batch scheduler
(\texttt{pack\_jobs.py}) packed GPU jobs on uicgpu to maximize throughput.
Tool use was unrestricted within the workspace.

\subsection{Scoring rubric}
We scored each paper on two axes, following the schema in
\texttt{scoring/SCHEMA.md} (reproduced in the Appendix):
\begin{itemize}
  \item \textbf{Coverage} (1--10): fraction of paper's contributions reproduced.
  10 = every major element at or near paper's scale; 8 = core method and main
  results, minor simplification; 6 = core method but significant scale/feature
  simplification; 4 = partial (e.g.\ comparator arm only); 2 = tangentially
  related demo; 1 = not replicated.
  \item \textbf{Agreement} (1--10): quantitative/qualitative fit to paper.
  10 = within 5--10\% everywhere; 8 = within 20\% with correct trends; 6 =
  qualitative match, 20--50\% metric gap; 4 = mechanism-level only; 2 = weak;
  1 = contradicts.
\end{itemize}
We emphasize that the two axes are intentionally orthogonal: a small
simplification of a paper's architecture (low coverage) can still reproduce the
main numerical finding exactly (high agreement), and a full-scale reimplementation
(high coverage) can diverge numerically due to force-field, stochastic-seed, or
hyperparameter choices (low agreement).


\clearpage
